% Chapter 5, section 4: Upgrading and monitoring
% Write in plain LaTeX. Code of a guide step: \guide{<step>}; commands: \begin{shell}...\end{shell};
% a screenshot: \screenshot{file}{caption}; a gap in the draft: \placeholder{...}.
% See book/latex/queen.sty for the full list.
%
% Draft by Claude, 2026-10-23. The Compose roll ran from 2.5.0 to 2.6.0 (17.6 s, every message
% kept); /health, the probe and the metrics ran on a 2.6.0 cluster with two nodes stopped for the
% probe example. The kubectl lines come from the site's Kubernetes and operator pages, not run.

\section{Upgrading and monitoring}

\subsection{Upgrade one node at a time}

Stop one node, start the new release on the same data, and wait until its \texttt{/health}
answers 200 before you touch the next. Healthy means the node knows the leader and has caught up
(it is less than 1,000 entries or two seconds behind), so it is back in the majority before the
next one stops.

In the Compose cluster, put the new version in \texttt{.env} and let \texttt{-{}-wait} hold each
step until the node is healthy:

\begin{shell}
echo "QUEEN_VERSION=<new-version>" >> .env
for n in 1 2 3; do docker compose up -d --wait queen-$n; done
\end{shell}

We rolled three nodes from 2.5.0 to 2.6.0 this way in 18 seconds. Every message pushed before the
roll was there after it, and pushing one of them again was answered \texttt{duplicate}: the
deduplication memory survived too.

On your own machines, do the same by hand, one machine at a time:

\begin{shell}
docker stop -t 60 queen && docker rm queen
# the same docker run as before, with the new version at the end, then:
until curl -sf -o /dev/null http://10.0.0.11:6632/health; do sleep 1; done
\end{shell}

On Kubernetes, change the image. The StatefulSet replaces one pod at a time, highest number first,
and waits for each to be Ready. The operator does the same, and restarts the leader last, so the
leadership moves only once:

\begin{shell}
kubectl -n queen set image statefulset/queen queen=ghcr.io/queen-mq/queen:<new-version>
kubectl -n queen patch queencluster queen --type merge \
  -p '{"spec":{"version":"<new-version>"}}'
\end{shell}

(The first line is for the manifest, the second for the operator.) A single node has no one to
take over: its upgrade is a stop and a start, a few seconds without service. Give the stop 60
seconds, because pops that are waiting for messages may take up to 30 seconds to finish.

\why{two nodes of three must be up for anything to commit. One node at a time, each one healthy
before the next, keeps two up at every moment, and the clients see at most a slower answer.}

\subsection{The way back}

A release can change how data is written. It does so only when every node runs it: until the last
node is upgraded, the cluster keeps writing the old format, and the new feature stays off. The
format is the \texttt{clusterVersion} in \texttt{/health}. While it has not moved, you can go back
by starting the old release on the same data. Once it moves, an older release refuses to start on
that data.

To keep the way back open while a new release proves itself, set
\texttt{QUEEN\_RAFT\_CLUSTER\_VERSION\_MS=0} on every node: the version then never moves. Remove the
setting when you decide to stay, and the version moves within a second. 2.5.0 and 2.6.0 write the
same format (version 7).

\why{a rolling upgrade means old and new nodes run side by side for a while. If a new node wrote
something an old node cannot read, the old node would stop. Waiting for the last node is what makes
the roll safe in both directions, up to the moment you choose.}

\queenfig{rolling-upgrade}{A rolling upgrade, and when the data format moves.}

\subsection{What to watch}

Every node serves three things on port 6632: the dashboard at \texttt{/}, its health at
\texttt{/health}, and its metrics at \texttt{/metrics/prometheus}. Each one describes the node that
answered, so probe and scrape every node at its own address, never through a load balancer.

\paragraph{The dashboard.} Its System page shows the replicated log and every node, with its role,
its term, how far it has applied the log and how much disk it uses. A bar at the top of every page
says when a node refuses writes because its disk is full.

\paragraph{/health.} It answers 200 \texttt{healthy} when the node knows a leader, is less than two
seconds behind, and the leader has heard from a majority in the last five seconds. Otherwise it
answers 503 \texttt{settling}.

\begin{shell}
curl -s http://10.0.0.13:6632/health
{"status":"healthy","engine":"raft","version":"2.6.0",
 "raft":{"role":"leader","leader":true,"term":1,"applied":5,"commit":5,"lag":0,
  "quorumAckMs":8,"storageReady":true,"clusterVersion":7,"kinds":7,
  "apply":{"failure":null,"skipped":[]}}}
\end{shell}

\begin{options}{Field}{Healthy}{What it means}
role & & \texttt{leader}, \texttt{follower}, \texttt{learner} (joining, not voting yet), \texttt{candidate} or \texttt{stopped}. \\
applied, commit & equal & How far this node has applied the log, and how far the cluster has committed it. \\
lag & 0 & Milliseconds this node is behind the leader. \\
quorumAckMs & small & Milliseconds since the leader last heard from a majority. It grows when the majority is lost. \\
apply.failure & null & The entry that stopped this node, if one did (see Recovery). \\
clusterVersion, kinds & & The format the cluster writes, and the newest this build can read. \\
\end{options}

Use \texttt{/health} for readiness. Use a plain TCP check on port 6632 for liveness, never
\texttt{/health}: a 503 during an election is no reason to restart a node.

\paragraph{A write probe.} \texttt{/health} reports what a node knows. It writes nothing, so it
cannot prove that the cluster commits. Three seconds after we stopped two nodes of three, the
third still answered 200, though no write could commit any more. So, beside
\texttt{/health}, write a small KV key through every node every minute, and alert when it fails or
takes more than a few seconds:

\queenfig{health-vs-probe}{After a lost majority: /health trails, the write probe does not.}

\begin{shell}
curl -s -m 10 -X PUT http://10.0.0.11:6632/api/v1/kv/ops/probe \
  -H 'content-type: application/json' -d '{"value":"probe","ttlSeconds":60}'
{"applied":true,"index":0,"key":"probe","op":"put","value":"probe","version":1}
\end{shell}

On a healthy cluster it answers in a few tens of milliseconds. The key expires after a minute, so the probe
leaves nothing behind.

\paragraph{Metrics.} These are the series worth an alert:

\begin{options}{Series}{Alert when}{What it means}
queen\_raft\_storage\_full & = 1 & The disk is 85\% full: writes that grow it answer 507. \\
queen\_raft\_index & committed $-$ applied keeps growing & The node applies the log slower than the cluster commits it. \\
queen\_raft\_admit\_total & \texttt{refused} rises & Writes are refused because the node is overloaded (429). \\
queen\_dlq\_depth\_by\_queue & rises & Messages are piling up in a dead-letter queue. \\
queen\_queue\_pop\_lag\_\allowbreak milliseconds & high & Consumers are behind. \\
queen\_uptime\_seconds & < 300 & The node restarted. \\
\end{options}

For memory, compare the process's own memory with its limit (\texttt{container\_memory\_rss} on
Kubernetes). The resident size also counts files the kernel keeps in its cache and drops before it
runs out of memory.

\paragraph{Log lines.} A few lines deserve an alert of their own:

\begin{options}{Line}{Level}{What it means}
FATAL: ... & ERROR & The node will not start. The line says why. \\
APPLY FAILED at entry N & ERROR & An entry stopped this node (see Recovery). \\
raft storage pressure changed full=true & WARN & The disk gate closed. \\
refused this node's QUEEN\_RAFT\_TOKEN & WARN & The nodes do not share one token. \\
the process exits to load a received snapshot & ERROR & Exit code 75. Normal once; trouble if it repeats. \\
\end{options}

\why{a cluster hides a failure well, which is the point of it, and also the danger: two nodes
serve while the third is broken, and nobody notices until a second one goes. The probe and these
alerts are how you notice while one failure is still free.}

On the site: \href{https://queenmq.com/manage/upgrades-monitoring/}{queenmq.com/manage/upgrades-monitoring}. More there: \href{https://queenmq.com/reference/metrics/}{every metric}.
